\documentclass[11pt,a4paper]{article}
\usepackage[utf8]{inputenc}
\usepackage[T1]{fontenc}
\usepackage[italian]{babel}
\usepackage{lmodern}
\usepackage{microtype}
\usepackage[margin=2cm,headheight=14pt]{geometry}
\usepackage{amsmath,amssymb,amsthm,mathtools,amsfonts,bm,siunitx}
\usepackage{booktabs,tabularx,enumitem,float}
\usepackage{xcolor}
\usepackage[most]{tcolorbox}
\usepackage{fancyhdr}
\usepackage{listings}
\usepackage{hyperref}
\usepackage{bookmark}

\definecolor{navyblue}{RGB}{10, 45, 90}
\definecolor{coolblue}{RGB}{28, 110, 164}
\definecolor{forestgreen}{RGB}{20, 115, 60}
\definecolor{warmamber}{RGB}{195, 110, 10}
\definecolor{crimsonred}{RGB}{175, 30, 45}
\definecolor{subtlegray}{RGB}{245, 247, 250}
\definecolor{codebg}{RGB}{240, 242, 245}

\hypersetup{
    colorlinks=true,
    linkcolor=navyblue,
    urlcolor=coolblue,
    citecolor=forestgreen
}

\pagestyle{fancy}
\fancyhf{}
\fancyhead[L]{\textbf{\textsf{\textcolor{navyblue}{Statistica I --- Soluzione Ufficiale Dettagliata}}}}
\fancyhead[R]{\textsf{\textcolor{coolblue}{II Appello (Luglio 2020)}}}
\fancyfoot[C]{\thepage}
\renewcommand{\headrulewidth}{0.6pt}

\newtcolorbox{problemabox}[1]{
    enhanced,
    breakable,
    colback=white,
    colframe=navyblue,
    coltitle=white,
    fonttitle=\bfseries\sffamily,
    title={#1},
    attach boxed title to top left={yshift=-2mm, xshift=3mm},
    boxed title style={colback=navyblue, boxrule=0pt, sharp corners, rounded corners=north},
    boxrule=1.2pt,
    sharp corners,
    rounded corners=south,
    top=4mm, bottom=3mm, left=3mm, right=3mm
}

\newtcolorbox{soluzionebox}[1][Svolgimento Dettagliato]{
    enhanced,
    breakable,
    colback=subtlegray,
    colframe=forestgreen!80!black,
    coltitle=white,
    fonttitle=\bfseries\sffamily,
    title={#1},
    attach boxed title to top left={yshift=-2mm, xshift=3mm},
    boxed title style={colback=forestgreen!80!black, boxrule=0pt, sharp corners, rounded corners=north},
    boxrule=1pt,
    sharp corners,
    rounded corners=south,
    top=4mm, bottom=3mm, left=3mm, right=3mm
}

\newtcolorbox{esamebox}[1][Consiglio per l'Esame \& Aspetti Chiave]{
    enhanced,
    breakable,
    colback=crimsonred!4!white,
    colframe=crimsonred,
    coltitle=crimsonred,
    fonttitle=\bfseries\sffamily,
    title={#1},
    attach boxed title to top left={yshift=-1.5mm, xshift=2mm},
    boxed title style={colback=white, boxrule=0.8pt, colframe=crimsonred},
    boxrule=0.8pt,
    leftrule=3.5pt,
    sharp corners,
    top=3.5mm, bottom=2.5mm, left=3mm, right=3mm
}

\lstset{
    backgroundcolor=\color{codebg},
    basicstyle=\ttfamily\small,
    breaklines=true,
    frame=single,
    rulecolor=\color{navyblue!40},
    keywordstyle=\color{navyblue}\bfseries,
    commentstyle=\color{forestgreen}\itshape,
    stringstyle=\color{warmamber},
    showstringspaces=false
}

\newcommand{\EE}{\mathbb{E}}
\newcommand{\Prob}{\mathbb{P}}
\newcommand{\Var}{\operatorname{Var}}
\newcommand{\Cov}{\operatorname{Cov}}
\newcommand{\dx}{\mathrm{d}x}

\begin{document}

\begin{center}
    {\LARGE\bfseries\color{navyblue} Università di Pisa --- Corso di Laurea in Ingegneria Gestionale}\\[0.4em]
    {\Large\bfseries Statistica I (A.A. 2019/2020) --- II Appello (Luglio 2020)}\\[0.3em]
    {\large\textsf{Soluzione Completa, Rigorosa e Commentata Passo-Passo}}
\end{center}
\vspace{0.5em}
\hrule height 1.2pt
\vspace{1.5em}

% ==============================================================================
% PROBLEMA 1
% ==============================================================================
\begin{problemabox}{Problema 1 (Testo Integrale)}
Siano $A,B$ due eventi. Si calcolino $\Prob(A)$ e $\Prob(B \mid A)$ sapendo che $\Prob(B)=\frac{3}{10}$, $\Prob(A \mid B)=\frac{9}{10}$ e $\Prob(A \mid B^c)=\frac{1}{10}$.
\end{problemabox}

\begin{soluzionebox}
\textbf{1. Formula delle Probabilità Totali per $\Prob(A)$:}
Gli eventi $B$ e $B^c$ costituiscono una partizione fondamentale dello spazio campionario $\Omega$.
Poiché $\Prob(B) = \frac{3}{10} = 0.3$, la probabilità del complementare è $\Prob(B^c) = 1 - \Prob(B) = 1 - \frac{3}{10} = \frac{7}{10} = 0.7$.
Applicando il Teorema delle Probabilità Totali:
\begin{align*}
\Prob(A) &= \Prob(A \mid B)\Prob(B) + \Prob(A \mid B^c)\Prob(B^c) \\
&= \left(\frac{9}{10}\right)\left(\frac{3}{10}\right) + \left(\frac{1}{10}\right)\left(\frac{7}{10}\right) = \frac{27}{100} + \frac{7}{100} = \frac{34}{100} = \bm{0.3400} = \frac{17}{50}.
\end{align*}

\textbf{2. Formula di Bayes per la probabilità a posteriori $\Prob(B \mid A)$:}
Applicando la definizione di probabilità condizionata (Teorema di Bayes):
\[
\Prob(B \mid A) = \frac{\Prob(A \cap B)}{\Prob(A)} = \frac{\Prob(A \mid B)\Prob(B)}{\Prob(A)} = \frac{\frac{9}{10} \cdot \frac{3}{10}}{\frac{34}{100}} = \frac{\frac{27}{100}}{\frac{34}{100}} = \frac{27}{34} \approx \bm{0.7941} \quad (\bm{79.41\%}).
\]
\end{soluzionebox}

% ==============================================================================
% PROBLEMA 2
% ==============================================================================
\begin{problemabox}{Problema 2 (Testo Integrale)}
Si calcoli approssimativamente l'integrale
\[
\int_{-5}^2 \frac{\log(2 + \sin(x))}{1 + e^x} \, \dx,
\]
implementando in R il metodo di Monte Carlo. Si indichi inoltre il comando R per disegnare il grafico della funzione integranda nell'intervallo di integrazione e si mostri il grafico prodotto esportandone l'immagine in pdf.
\end{problemabox}

\begin{soluzionebox}
\textbf{1. Formulazione Monte Carlo:}
L'integrale $I = \int_{-5}^2 g(x) \, \dx$ con $g(x) = \frac{\log(2+\sin(x))}{1+e^x}$ ha estremi $a = -5, b = 2$ ($b - a = 7$).
Possiamo riscriverlo come:
\[
I = 7 \int_{-5}^2 g(x) \frac{1}{7} \, \dx = 7 \, \EE[g(X)], \qquad X \sim \operatorname{Unif}(-5, 2).
\]
Il valore numerico deterministico dell'integrale è $I \approx \bm{3.0181}$.

\textbf{2. Codice R per calcolo Monte Carlo e generazione PDF del grafico:}
\begin{lstlisting}[language=R]
# 1. Definizione della funzione integranda
g <- function(x) {
  log(2 + sin(x)) / (1 + exp(x))
}

# 2. Integrazione Monte Carlo
set.seed(123)
N <- 1000000
x_samples <- runif(N, min = -5, max = 2)
g_values <- g(x_samples)
I_hat <- 7 * mean(g_values)
se_hat <- 7 * sd(g_values) / sqrt(N)

cat(sprintf("Stima Monte Carlo: %.4f (SE: %.4f)\n", I_hat, se_hat))
# Output atteso: 3.0181

# 3. Tracciamento del grafico ed esportazione in PDF
pdf("grafico_funzione_integranda.pdf", width = 7, height = 5)
curve(g(x), from = -5, to = 2, col = "navyblue", lwd = 2,
      main = expression(paste("Grafico di ", g(x) == frac(log(2+sin(x)), 1+e^x))),
      xlab = "x", ylab = "g(x)", las = 1)
grid()
dev.off()
\end{lstlisting}
\end{soluzionebox}

% ==============================================================================
% PROBLEMA 3
% ==============================================================================
\begin{problemabox}{Problema 3 (Testo Integrale)}
Sia $\overline{X}_n$ la media empirica di $n$ copie indipendenti di una variabile $X$ di Poisson di parametro $\lambda=10$ e sia $Z_n$ la variabile ottenuta standardizzando $\overline{X}_n$. Si disegni a mano il grafico della densità di $Z_n$ per $n$ molto grande. Si calcoli inoltre approssimativamente usando R la probabilità che $Z_n$ sia maggiore di $2$ (sempre assumendo che $n$ sia molto grande).
\end{problemabox}

\begin{soluzionebox}
\textbf{1. Media, varianza e standardizzazione della media empirica:}
Poiché $X \sim \operatorname{Poiss}(\lambda=10)$, abbiamo $\EE[X] = 10$ e $\Var(X) = 10$.
Per la media empirica $\overline{X}_n = \frac{1}{n}\sum_{i=1}^n X_i$:
\[
\EE[\overline{X}_n] = 10, \qquad \Var(\overline{X}_n) = \frac{\Var(X)}{n} = \frac{10}{n} \implies \sigma_{\overline{X}_n} = \sqrt{\frac{10}{n}}.
\]
La variabile aleatoria standardizzata è:
\[
Z_n = \frac{\overline{X}_n - \EE[\overline{X}_n]}{\sigma_{\overline{X}_n}} = \frac{\overline{X}_n - 10}{\sqrt{10/n}}.
\]

\textbf{2. Comportamento asintotico per $n \to \infty$:}
In virtù del Teorema del Limite Centrale (TLC), per $n \to \infty$:
\[
Z_n \xrightarrow{\mathcal{D}} Z \sim \mathcal{N}(0, 1).
\]
La densità limite è la curva a campana di Gauss standard:
\[
\phi(z) = \frac{1}{\sqrt{2\pi}} e^{-z^2/2}.
\]
Il grafico presenta asse di simmetria in $z = 0$, punti di flesso in $z = \pm 1$, e valore massimo pari a $\phi(0) = \frac{1}{\sqrt{2\pi}} \approx 0.3989$.

\textbf{3. Calcolo di $\Prob(Z_n > 2)$ con R:}
Per $n$ grande:
\[
\Prob(Z_n > 2) \approx \Prob(Z > 2) = 1 - \Phi(2) = 1 - 0.97725 = \bm{0.0228} \quad (\bm{2.28\%}).
\]
Comando R:
\begin{lstlisting}[language=R]
# Probabilita che Zn > 2
p_val <- 1 - pnorm(2, mean = 0, sd = 1)
cat(sprintf("P(Zn > 2) = %.4f\n", p_val)) # 0.0228
\end{lstlisting}
\end{soluzionebox}

% ==============================================================================
% PROBLEMA 4
% ==============================================================================
\begin{problemabox}{Problema 4 (Testo Integrale)}
Si consideri una popolazione di legge $\operatorname{Unif}(-\theta, 4\theta)$ e si proponga per $\theta$ uno stimatore puntuale $\widehat{\theta}$ basato su un campione casuale $X_1, \dots, X_{2000}$.
Si scelga un file a caso e si dia una stima per $\theta$ usando $\widehat{\theta}$ e i dati del file prescelto.
\end{problemabox}

\begin{soluzionebox}
\textbf{1. Metodo dei Momenti per popolazione uniforme:}
Sia $X \sim \operatorname{Unif}(-\theta, 4\theta)$ con $\theta > 0$.
L'ampiezza dell'intervallo è $4\theta - (-\theta) = 5\theta$. Il valore atteso è il punto medio:
\[
\EE[X] = \frac{-\theta + 4\theta}{2} = \frac{3}{2}\theta.
\]
Ponendo il valore atteso teorico uguale alla media campionaria $\overline{X}_n = \frac{1}{n}\sum_{i=1}^n X_i$:
\[
\frac{3}{2}\theta = \overline{X}_n \implies \bm{\widehat{\theta}_{\text{MM}} = \frac{2}{3}\overline{X}_n}.
\]
Verifica di correttezza: $\EE[\widehat{\theta}_{\text{MM}}] = \frac{2}{3}\EE[\overline{X}_n] = \frac{2}{3} \cdot \frac{3}{2}\theta = \theta$ (stimatore corretto).

\textbf{2. Script R per caricare un dataset casuale e calcolare la stima:}
\begin{lstlisting}[language=R]
# Selezione casuale di un file tra i 200 disponibili
set.seed(42)
file_id <- sample(1:200, 1)
file_name <- sprintf("campione_%d.csv", file_id)

# Caricamento dati
dati <- read.csv(file_name)
x <- dati[, 1]
n <- length(x) # n = 2000

# Calcolo media campionaria e stima puntuale
x_bar <- mean(x)
theta_hat <- (2/3) * x_bar

cat(sprintf("File selezionato: %s\n", file_name))
cat(sprintf("Media campionaria: x_bar = %.4f\n", x_bar))
cat(sprintf("Stima puntuale theta_hat: %.4f\n", theta_hat))
\end{lstlisting}
\end{soluzionebox}

% ==============================================================================
% PROBLEMA 5
% ==============================================================================
\begin{problemabox}{Problema 5 (Testo Integrale)}
Siano $X,Y$ due variabili aleatorie indipendenti entrambe distribuite secondo una legge binomiale di parametri $n=1$ e $p=\frac{1}{3}$. Consideriamo le variabili aleatorie $Z=XY$ e $U=X-Y$.
\begin{enumerate}[label=(\roman*)]
    \item Si calcoli la funzione di massa di $Z$. Che valore atteso e che varianza ha $Z$?
    \item Qual è la probabilità che $Z \le X$?
    \item Quale variabile aleatoria si ottiene standardizzando $U$?
    \item Si calcoli la funzione di massa di $U$.
    \item Si calcoli la funzione di massa congiunta della coppia $(Z,U)$. Le variabili $Z,U$ sono indipendenti? Sono scorrelate?
\end{enumerate}
\end{problemabox}

\begin{soluzionebox}
\textbf{1. Distribuzione congiunta di $(X,Y)$:}
$X, Y \overset{\text{i.i.d.}}{\sim} \operatorname{Bernoulli}(1/3)$ con supporto $\{0, 1\}$:
\[
\Prob(X=1) = \Prob(Y=1) = \frac{1}{3}, \qquad \Prob(X=0) = \Prob(Y=0) = \frac{2}{3}.
\]
Le 4 configurazioni congiunte per indipendenza sono:
\[
\Prob(X=0, Y=0) = \frac{4}{9}, \quad \Prob(X=0, Y=1) = \frac{2}{9}, \quad \Prob(X=1, Y=0) = \frac{2}{9}, \quad \Prob(X=1, Y=1) = \frac{1}{9}.
\]

\textbf{2. Risoluzione dei singoli quesiti:}
\begin{enumerate}[label=(\roman*)]
    \item \textbf{Funzione di massa, attesa e varianza di $Z = XY$:}
    Il prodotto $Z$ assume valore $1$ se e solo se $X=1$ e $Y=1$, altrimenti vale $0$:
    \[
    p_Z(1) = \Prob(X=1, Y=1) = \frac{1}{3} \cdot \frac{1}{3} = \bm{\frac{1}{9}}, \qquad p_Z(0) = 1 - \frac{1}{9} = \bm{\frac{8}{9}}.
    \]
    Quindi $Z \sim \operatorname{Bernoulli}(1/9)$.
    \[
    \EE[Z] = \bm{\frac{1}{9}} \approx \bm{0.1111}, \qquad \Var(Z) = \frac{1}{9} \cdot \frac{8}{9} = \bm{\frac{8}{81}} \approx \bm{0.0988}.
    \]

    \item \textbf{Probabilità che $Z \le X$:}
    Poiché $Y \in \{0, 1\}$, abbiamo sempre $XY \le X$ per ogni realizzazione (se $X=0 \implies Z=0 \le 0$; se $X=1 \implies Z=Y \le 1$). L'evento $\{Z \le X\}$ è l'evento certo:
    \[
    \Prob(Z \le X) = \bm{1}.
    \]

    \item \textbf{Standardizzazione di $U = X - Y$:}
    $\EE[U] = \EE[X] - \EE[Y] = \frac{1}{3} - \frac{1}{3} = 0$.
    Per indipendenza: $\Var(U) = \Var(X) + \Var(Y) = \frac{2}{9} + \frac{2}{9} = \frac{4}{9} \implies \sigma_U = \frac{2}{3}$.
    La variabile standardizzata $U^*$ è:
    \[
    \bm{U^* = \frac{U - \EE[U]}{\sigma_U} = \frac{U}{2/3} = \frac{3}{2}U = \frac{3}{2}(X - Y)}.
    \]

    \item \textbf{Funzione di massa di $U = X - Y$:}
    Il supporto di $U$ è $\{-1, 0, 1\}$:
    \begin{align*}
    p_U(-1) &= \Prob(X=0, Y=1) = \bm{\frac{2}{9}}, \\
    p_U(0) &= \Prob(X=0, Y=0) + \Prob(X=1, Y=1) = \frac{4}{9} + \frac{1}{9} = \bm{\frac{5}{9}}, \\
    p_U(1) &= \Prob(X=1, Y=0) = \bm{\frac{2}{9}}.
    \end{align*}

    \item \textbf{Funzione di massa congiunta di $(Z, U)$, Indipendenza e Correlazione:}
    Costruiamo la tabella della PMF congiunta $p_{Z,U}(z, u)$:
    \begin{center}
    \begin{tabular}{c|ccc|c}
    \toprule
    $p_{Z,U}(z, u)$ & $u = -1$ & $u = 0$ & $u = 1$ & $p_Z(z)$ \\
    \midrule
    $z = 0$ & $2/9$ & $4/9$ & $2/9$ & $8/9$ \\
    $z = 1$ & $0$ & $1/9$ & $0$ & $1/9$ \\
    \midrule
    $p_U(u)$ & $2/9$ & $5/9$ & $2/9$ & $1$ \\
    \bottomrule
    \end{tabular}
    \end{center}
    \begin{itemize}
        \item \emph{Indipendenza:} Notiamo che $p_{Z,U}(1, 1) = 0$, mentre $p_Z(1) p_U(1) = \frac{1}{9} \cdot \frac{2}{9} = \frac{2}{81} \neq 0$.
        Poiché $p_{Z,U}(1,1) \neq p_Z(1)p_U(1)$, le variabili $Z$ e $U$ \textbf{NON sono indipendenti}.
        \item \emph{Correlazione:} Calcoliamo $\EE[ZU]$:
        \[
        \EE[ZU] = 1 \cdot (-1) \cdot 0 + 1 \cdot 0 \cdot \frac{1}{9} + 1 \cdot 1 \cdot 0 = 0.
        \]
        Dato che $\EE[U] = 0$, abbiamo:
        \[
        \Cov(Z, U) = \EE[ZU] - \EE[Z]\EE[U] = 0 - \left(\frac{1}{9}\right)(0) = \bm{0}.
        \]
        Essendo la covarianza nulla, $Z$ e $U$ sono \textbf{scorrelate} ($\rho_{ZU}=0$), fornendo un perfetto esempio di variabili dipendenti ma con indice di correlazione nullo!
    \end{itemize}
\end{enumerate}
\end{soluzionebox}

% ==============================================================================
% PROBLEMA 6
% ==============================================================================
\begin{problemabox}{Problema 6 (Testo Integrale)}
Per combattere un'influenza viene introdotto un nuovo tipo di vaccino obbligatorio per tutti. Su $340$ individui estratti a caso dopo un anno, si scopre che $293$ non hanno contratto l'influenza, mentre i restanti $47$ hanno contratto l'influenza nonostante la vaccinazione. Si determini un valore $a_0$ tale che con confidenza del $95\%$ si possa affermare che, rispetto a tutta la popolazione vaccinata, la proporzione di individui non colpiti dall'influenza sia almeno $a_0$.
\end{problemabox}

\begin{soluzionebox}
\textbf{1. Modello Bernoulliano e Quantità Pivotale Asintotica:}
Sia $p$ la reale proporzione (incognita) di individui non colpiti nella popolazione vaccinata.
Il campione è di dimensione $n = 340$, con numero di successi $k = 293$.
La proporzione campionaria stimata è:
\[
\widehat{p} = \frac{k}{n} = \frac{293}{340} \approx \bm{0.8618} \quad (\bm{86.18\%}).
\]
L'errore standard asintotico stimato è:
\[
\operatorname{SE}(\widehat{p}) = \sqrt{\frac{\widehat{p}(1-\widehat{p})}{n}} = \sqrt{\frac{0.86176 \cdot 0.13824}{340}} = \sqrt{\frac{0.11913}{340}} \approx \bm{0.01872}.
\]

\textbf{2. Intervallo di Confidenza Unilaterale Inferiore al 95\%:}
Cerchiamo una soglia minima $a_0$ tale che:
\[
\Prob(p \ge a_0) = 1 - \alpha = 0.95.
\]
Sfruttando l'approssimazione normale standard per grandi campioni:
\[
\frac{\widehat{p} - p}{\sqrt{\frac{\widehat{p}(1-\widehat{p})}{n}}} \approx \mathcal{N}(0, 1).
\]
La condizione unilaterale corrisponde al quantile critico $z_{0.05} = 1.6449$ (dove $\Phi(1.6449) = 0.95$):
\[
a_0 = \widehat{p} - z_{0.05} \cdot \operatorname{SE}(\widehat{p}) = 0.86176 - 1.6449 \cdot 0.01872 = 0.86176 - 0.03079 = \bm{0.8310} \quad (\bm{83.10\%}).
\]

\textbf{Conclusione:}
Con un livello di confidenza del $95\%$, possiamo affermare che la proporzione di individui protetti dal vaccino nella popolazione complessiva è pari ad \textbf{almeno l'83.10\%} ($a_0 = \bm{0.8310}$).
\end{soluzionebox}

\end{document}
